% Two Koriyama slides rebuilt in Beamer (example) with infoblocks.sty, to compare with
% the Quarto deck (../quarto/, slides 6 and 8). Content as in the Quarto deck
% (Kakushokai wording).
\documentclass[aspectratio=169,10pt]{beamer}

\usepackage{tabularray}
\usepackage{fontspec}
\usefonttheme{professionalfonts}
\usepackage[math-style=TeX]{unicode-math}
\setmainfont[Ligatures=TeX]{Noto Sans}
\setmathfont[math-style=ISO,bold-style=ISO,vargreek-shape=TeX]{Lete Sans Math}
\usepackage{zxjatype}
\usepackage{xeCJK}
\setCJKmainfont{Noto Sans CJK JP}
\setCJKsansfont{Noto Sans CJK JP}

\usetheme{JECS}
\usepackage{infoblocks}

\title{infoblocks demo}
\date{2026年11月16日}

\begin{document}

% ==========================================================================
\begin{frame}[t]{ロードマップ：調査の構成}
	\begin{columns}[T,totalwidth=\textwidth]
		\begin{column}{0.325\textwidth}
			\begin{infocard}[navy][equal height group=kosei]{users}{全体調査}
				\stat{約10万人}
				\begin{itemize}
					\item 質問票調査（保護者）：妊娠中〜小学生まで年2回、中学生以降はweb
					\item 質問票調査（本人）：小学4年生〜6年生まで年1回、中学生以降はweb
					\item 学童期検査（小2・小6）、乳歯調査（小4）
					\item 青年期調査（高1）
				\end{itemize}
			\end{infocard}
		\end{column}
		\begin{column}{0.325\textwidth}
			\begin{infocard}[teal][equal height group=kosei]{magnifying-glass}{詳細調査}
				\stat[teal]{約5,000人}
				{\tiny\color{black!55}2013--14年度生まれから抽出\par}
				\begin{itemize}
					\item 訪問調査（1.5歳・3歳）：ダニアレルゲン、VOCs、ハウスダストなど
					\item 医学的検査・精神神経発達検査（2歳・4歳・6歳〈医学のみ〉・8歳・10歳・12歳）
					\item 新版K式発達検査、CAT、10歳でWISC
				\end{itemize}
			\end{infocard}
		\end{column}
		\begin{column}{0.325\textwidth}
			\begin{infocard}[red][equal height group=kosei]{vault}{生体試料の保管・分析}
				\stat[red]{500万本超}
				\begin{itemize}
					\item 血液、臍帯血、母乳、尿、毛髪、乳歯を長期保管
					\item 新たな測定法の確立後に、遡及的な測定が可能
					\item 化学分析、遺伝子解析
				\end{itemize}
			\end{infocard}
		\end{column}
	\end{columns}

	\vspace{4pt}
	{\tiny\color{black!55}成人に達する年度以降も、データクリーニング、解析、論文化、発表（論文、学会）を継続する。}
\end{frame}

% ==========================================================================
\begin{frame}[t]{13歳以降調査に関する継続の意思表示状況}
	\begin{columns}[T,totalwidth=\textwidth]
		\begin{column}{0.325\textwidth}
			\stattile{users}{96.1\%}{親権者・継続者総数}
		\end{column}
		\begin{column}{0.325\textwidth}
			\stattile[red]{thumbs-up}{51.2\%}{親権者・積極的意思表示}
		\end{column}
		\begin{column}{0.325\textwidth}
			\stattile[teal]{laptop}{41.2\%}{本人（子ども）・ログイン数}
		\end{column}
	\end{columns}

	\vspace{3pt}
	{\tiny\color{black!55}合計（2011--13年度生まれ、75,420人）。2026年8月24日現在。}

	\vspace{6pt}
	\begin{tblr}{width=\textwidth, colspec={Q[l,m,wd=2.2cm] *{4}{X[l,m]}},
		colsep=4pt, rowsep=3pt, cells={font=\scriptsize},
		row{1}={font=\scriptsize\bfseries, fg=ibnavy},
		cell{1}{5}={fg=ibred}, column{1}={font=\scriptsize\bfseries}}
		& 2011年度生まれ\par{\tiny\mdseries\color{black!55}15,497人・小6は2023年度}
		& 2012年度生まれ\par{\tiny\mdseries\color{black!55}29,098人・小6は2024年度}
		& 2013年度生まれ\par{\tiny\mdseries\color{black!55}35,610人・小6は2025年度}
		& 2014年度生まれ\par{\tiny\mdseries\color{black!55}18,097人・小6は2026年度} \\
		{\color{ibnavy}継続者総数}
		& \meter{86.0}[2.4cm] & \meter{98.7}[2.4cm] & \meter{98.5}[2.4cm] & \meter{99.4}[2.4cm] \\
		{\color{ibred}積極的意思表示}
		& \meter[red]{51.8}[2.4cm] & \meter[red]{50.4}[2.4cm] & \meter[red]{51.7}[2.4cm] & \meter[red]{38.6}[2.4cm] \\
		{\color{ibtealink}本人ログイン}
		& \meter[teal]{41.2}[2.4cm] & \meter[teal]{40.3}[2.4cm] & \meter[teal]{42.1}[2.4cm] & \meter[teal]{25.2}[2.4cm] \\
	\end{tblr}
\end{frame}

\end{document}
